\chapter{Clase 20: Elementos orbitales en el problema de dos cuerpos}

Esta clase cierra el puente entre la solución analítica del problema de dos cuerpos y su uso práctico en astronomía y astrodinámica. Después de haber formulado la órbita en el sistema natural y de haber introducido las rotaciones necesarias para orientarla en el espacio, el siguiente paso es condensar toda esa información en un conjunto mínimo de parámetros geométricos y dinámicos: los \textbf{elementos orbitales}.

\section{Repaso: sistema natural y rotación plana}
La ecuación de la trayectoria relativa en coordenadas polares se expresa de forma canónica en el \textbf{sistema natural de la cónica}, también llamado sistema perifocal, cuyo eje \(x'\) coincide con la dirección del vector de Laplace \(\vec{e}\) y cuyo origen se sitúa en el foco gravitacional:
\[
 r = \frac{p}{1 + e\cos f}
\]
donde \(p = h^2/\mu\) es el semilatus rectum, \(e = |\vec{e}|\) la excentricidad y \(f\) la verdadera anomalía.

En este mismo sistema, las coordenadas cartesianas se obtienen proyectando \(\vec{r}\) sobre la base rotada \(\{\hat{i}', \hat{j}'\}\):
\[
 x' = r\cos f, \quad y' = r\sin f
\]

El \textbf{sistema del observador} \(\{\hat{i}, \hat{j}, \hat{k}\}\) rara vez coincide con el sistema natural. En el caso plano, la orientación relativa se describe mediante una rotación dada por
\[
\begin{pmatrix} x \\ y \end{pmatrix} =
\begin{pmatrix}
\cos\omega & -\sin\omega \\
\sin\omega & \cos\omega
\end{pmatrix}
\begin{pmatrix} x' \\ y' \end{pmatrix}
\]
Esta idea es la base para extender el análisis al espacio tridimensional.

\begin{importante}
Los elementos orbitales surgen precisamente como una manera compacta de describir la forma, la orientación y la posición instantánea del cuerpo dentro de su órbita kepleriana.
\end{importante}

\section{La pesadilla de los cuadrantes}
En la determinación analítica de ángulos orbitales, el uso directo de funciones trigonométricas inversas como \(\arccos\) introduce una ambigüedad fundamental: su rango principal está restringido a \([0,\pi]\). Esto impide distinguir entre configuraciones geométricas simétricas respecto al eje de referencia.

\subsection{Ambigüedad en la longitud del nodo ascendente \texorpdfstring{$\Omega$}{Omega}}
El vector nodo \(\vec{n}\) se define como la intersección del plano orbital con el plano de referencia, proyectada sobre dicho plano:
\[
\vec{n} = \hat{k} \times \vec{h}
\]

Si calculamos
\[
\Omega = \arccos\left(\frac{n_x}{n}\right),
\]
obtenemos un valor en \([0,\pi]\). Sin embargo, si \(n_y < 0\), el nodo ascendente se encuentra en el semiplano inferior, por lo que el ángulo real debe corregirse:
\[
\Omega =
\begin{cases}
\arccos(n_x/n), & \text{si } n_y \ge 0 \\
2\pi - \arccos(n_x/n), & \text{si } n_y < 0
\end{cases}
\]

\subsection{Ambigüedad en el argumento del periapsis \texorpdfstring{$\omega$}{omega}}
De manera análoga,
\[
\omega = \arccos\left(\frac{\vec{n}\cdot\vec{e}}{ne}\right)
\]
devuelve un valor principal. La corrección de cuadrante depende del signo de la componente \(z\) del vector de excentricidad, o equivalentemente del sentido relativo entre \(\vec{n}\), \(\vec{e}\) y \(\vec{h}\):
\[
\omega =
\begin{cases}
\arccos\left(\frac{\vec{n}\cdot\vec{e}}{ne}\right), & \text{si } e_z \ge 0 \\
2\pi - \arccos\left(\frac{\vec{n}\cdot\vec{e}}{ne}\right), & \text{si } e_z < 0
\end{cases}
\]

\begin{importante}
La determinación unívoca de ángulos orbitales requiere usar \(\operatorname{atan2}(y,x)\) o, de forma equivalente, corregir los cuadrantes inspeccionando el signo de las componentes transversales adecuadas.
\end{importante}

\section{Extensión a tres dimensiones: la cónica levantada}
Cuando el plano orbital no coincide con el plano de referencia, hacen falta tres rotaciones sucesivas para alinear el sistema perifocal con el sistema del observador. Esta secuencia corresponde a los \textbf{ángulos de Euler} en la convención \(3\)-\(1\)-\(3\), ampliamente usada en mecánica celeste.

\subsection{Línea de nodos y referencia planetaria}
Introducimos tres objetos geométricos clave:
\begin{itemize}
    \item el \textbf{plano de referencia}, usualmente el plano ecuatorial terrestre o el plano de la eclíptica,
    \item la \textbf{línea de nodos}, intersección entre el plano orbital y el plano de referencia,
    \item el \textbf{nodo ascendente}, punto donde el cuerpo cruza el plano de referencia moviéndose hacia el hemisferio norte \((z>0)\).
\end{itemize}

\subsection{Definición de los ángulos de orientación}
Los ángulos espaciales fundamentales son:
\begin{enumerate}
    \item \textbf{Longitud del nodo ascendente \(\Omega\):} ángulo en el plano de referencia desde el eje \(x\) de referencia hasta la línea de nodos.
    \item \textbf{Inclinación orbital \(I\):} ángulo entre el vector momento angular \(\vec{h}\) y el eje \(z\) de referencia.
    \item \textbf{Argumento del periapsis \(\omega\):} ángulo medido dentro del plano orbital desde la línea de nodos hasta la dirección del periapsis \((\vec{e})\).
\end{enumerate}

Estos tres ángulos, junto con \(p\), \(e\) y \(f\), determinan completamente la orientación y forma de la órbita en el espacio tridimensional.

\section{Matriz de rotación de Euler en 3D}
La transformación desde el sistema perifocal al sistema del observador se realiza mediante el producto de tres matrices de rotación elementales. La convención estándar aplica las rotaciones en el orden:
\begin{enumerate}
    \item \(\mathbf{R}_z(\omega)\): alinea el eje perifocal \(x'\) con la línea de nodos,
    \item \(\mathbf{R}_x(I)\): inclina el plano orbital respecto al plano de referencia,
    \item \(\mathbf{R}_z(\Omega)\): alinea la línea de nodos con el eje \(x\) del observador.
\end{enumerate}

La matriz de rotación de Euler resultante es:
\[
\mathbf{R}_{\text{Euler}} = \mathbf{R}_z(\Omega)\,\mathbf{R}_x(I)\,\mathbf{R}_z(\omega)
\]
Desarrollando el producto matricial:
\[
\mathbf{R}_{\text{Euler}} =
\begin{pmatrix}
\cos\Omega\cos\omega - \sin\Omega\sin\omega\cos I & -\cos\Omega\sin\omega - \sin\Omega\cos\omega\cos I & \sin\Omega\sin I \\
\sin\Omega\cos\omega + \cos\Omega\sin\omega\cos I & -\sin\Omega\sin\omega + \cos\Omega\cos\omega\cos I & -\cos\Omega\sin I \\
\sin\omega\sin I & \cos\omega\sin I & \cos I
\end{pmatrix}
\]

La posición en el sistema del observador se obtiene como
\[
\vec{r}_{\text{obs}} = \mathbf{R}_{\text{Euler}}\,\vec{r}_{\text{per}}
\]
donde
\[
\vec{r}_{\text{per}} = (r\cos f,\; r\sin f,\; 0)^T.
\]

La transformación inversa usa la transpuesta, válida por la ortogonalidad de \(\mathbf{R}_{\text{Euler}}\):
\[
\vec{r}_{\text{per}} = \mathbf{R}_{\text{Euler}}^T\,\vec{r}_{\text{obs}}
\]

\begin{importante}
La matriz de Euler resume de manera compacta toda la orientación espacial de la órbita. La forma de la cónica se mantiene intacta; lo que cambia es su orientación con respecto al observador.
\end{importante}

\section{Elementos orbitales clásicos y sus variantes}
Los seis parámetros que describen completamente una órbita kepleriana se agrupan en los \textbf{elementos orbitales clásicos}:
\begin{center}
\begin{tabular}{lll}
\textbf{Elemento} & \textbf{Símbolo} & \textbf{Significado geométrico} \\
Semilatus rectum & \(p\) & tamaño o escala de la cónica \\
Excentricidad & \(e\) & forma de la cónica \\
Inclinación & \(I\) & orientación del plano orbital \\
Longitud del nodo ascendente & \(\Omega\) & orientación de la línea de nodos \\
Argumento del periapsis & \(\omega\) & orientación del periapsis en el plano \\
Verdadera anomalía & \(f\) & posición instantánea del cuerpo
\end{tabular}
\end{center}

\subsection{Variantes descriptivas}
Dependiendo del tipo de problema, el parámetro \(p\) puede reemplazarse por otras cantidades de tamaño:
\begin{itemize}
    \item \textbf{Elementos asteroidales:} \((a, e, I, \Omega, \omega, f)\), donde \(a\) es el semieje mayor. Son útiles para órbitas cerradas \((e<1)\), con
\[
 p = a(1-e^2).
\]
    \item \textbf{Elementos cometarios:} \((q, e, I, \Omega, \omega, f)\), donde \(q\) es la distancia al periapsis. Son ventajosos para órbitas casi parabólicas \((e\approx 1)\), con
\[
 p = q(1+e).
\]
\end{itemize}

\section{Determinación de órbita: del vector de estado a los elementos}
El vector de estado \(\{\vec{r},\vec{v}\}\) contiene toda la información necesaria para recuperar los elementos orbitales. El procedimiento analítico estándar es el siguiente.

\subsection{Paso 1. Vectores fundamentales}
\[
\vec{h} = \vec{r} \times \vec{v}, \qquad
\vec{e} = \frac{1}{\mu}(\vec{v} \times \vec{h}) - \frac{\vec{r}}{r}, \qquad
\vec{n} = \hat{k} \times \vec{h}
\]

\subsection{Paso 2. Parámetros de tamaño y forma}
\[
 h = |\vec{h}|, \qquad e = |\vec{e}|, \qquad p = \frac{h^2}{\mu}
\]

\subsection{Paso 3. Orientación del plano orbital}
\[
 I = \arccos\left(\frac{h_z}{h}\right), \qquad n = |\vec{n}|
\]
\[
\Omega =
\begin{cases}
\arccos(n_x/n), & n_y \ge 0 \\
2\pi - \arccos(n_x/n), & n_y < 0
\end{cases}
\]

\subsection{Paso 4. Orientación del periapsis y posición instantánea}
\[
\omega =
\begin{cases}
\arccos\left(\frac{\vec{n}\cdot\vec{e}}{ne}\right), & e_z \ge 0 \\
2\pi - \arccos\left(\frac{\vec{n}\cdot\vec{e}}{ne}\right), & e_z < 0
\end{cases}
\]
\[
 f =
\begin{cases}
\arccos\left(\frac{\vec{e}\cdot\vec{r}}{er}\right), & \vec{r}\cdot\vec{v} \ge 0 \\
2\pi - \arccos\left(\frac{\vec{e}\cdot\vec{r}}{er}\right), & \vec{r}\cdot\vec{v} < 0
\end{cases}
\]

Este algoritmo permite pasar de \((x,y,z,v_x,v_y,v_z)\) a \((p,e,I,\Omega,\omega,f)\) de manera biunívoca, siempre que se apliquen correctamente las correcciones de cuadrante.

\begin{importante}
El vector de estado y los elementos orbitales son dos maneras equivalentes de describir la misma órbita kepleriana. Uno privilegia la formulación dinámica instantánea; el otro, la interpretación geométrica global.
\end{importante}

\section{Órbita osculatriz y elementos osculantes}
En un sistema real sometido a perturbaciones no keplerianas —por ejemplo, achatamiento terrestre, presión de radiación o terceros cuerpos— la trayectoria no es una cónica exacta. Sin embargo, en cualquier instante \(t_0\), puede construirse la \textbf{órbita osculatriz}: la cónica kepleriana que sería seguida por el cuerpo si, a partir de \(t_0\), desaparecieran todas las perturbaciones.

Sus propiedades principales son:
\begin{itemize}
    \item la órbita osculatriz es \textbf{tangente} a la trayectoria real en \(t_0\),
    \item comparte el mismo vector de estado \(\{\vec{r}(t_0),\vec{v}(t_0)\}\),
    \item sus elementos \((p(t), e(t), I(t), \Omega(t), \omega(t), f(t))\) dependen del tiempo y se denominan \textbf{elementos osculantes}.
\end{itemize}

La evolución temporal de estos elementos se describe mediante ecuaciones de variación de parámetros, como las ecuaciones de Lagrange o de Gauss, que relacionan las fuerzas perturbadoras con las derivadas temporales de los elementos orbitales.

\section{Síntesis final}
El desarrollo de esta clase cierra el puente entre la integración analítica del problema de dos cuerpos y su aplicación observacional:
\begin{enumerate}
    \item La geometría de la cónica se parametriza eficientemente mediante seis elementos orbitales.
    \item La orientación en el espacio se resuelve mediante tres rotaciones de Euler, cuya matriz conecta el sistema perifocal con el sistema del observador.
    \item La ambigüedad trigonométrica se resuelve mediante correcciones de cuadrante o mediante funciones trigonométricas con información de signo, como \(\operatorname{atan2}\).
    \item El vector de estado \(\{\vec{r},\vec{v}\}\) es suficiente y necesario para recuperar unívocamente los elementos orbitales.
    \item El concepto de órbita osculatriz permite extender el formalismo kepleriano a sistemas perturbados, sentando las bases para la teoría de perturbaciones y la propagación numérica de órbitas.
\end{enumerate}

\begin{importante}
Este marco constituye una de las columnas vertebrales de la astrodinámica moderna: conecta geometría, dinámica, observación y cómputo en un único lenguaje orbital coherente.
\end{importante}